% Response fragment: requires xcolor and natbib; use AE1_references.bib.
% The last paragraph follows the author's confirmation that Sections 1.2 and 3.3
% have been revised. Reconcile the supplied manuscript additions before submission.
{\color{blue}
\noindent\textbf{Response:}
We thank the Associate Editor for suggesting this perspective.
We agree that the value of a category-specific reference lies both in
the modeling knowledge it conveys and in the procedures that guide
its application. We have clarified these complementary roles and
their connection to recent work on agent skills.

Prior work distinguishes declarative knowledge, concerning facts
and concepts, from procedural knowledge, concerning how to perform
a task \citep{ae1_li2024metacognitive}. In LEAN-LLM-OPT, Ref-Data
provides external, case-based modeling knowledge through reference
problem descriptions, associated data, expert formulations, and
annotations identifying relevant data. These resources illustrate
model structures and how data enter a formulation. A worked example
can also convey procedural knowledge when organized as a modeling
demonstration; we therefore distinguish the roles of knowledge and
procedure without treating them as disjoint resources.

Retrieval-augmented generation (RAG) describes a mechanism for
conditioning generation on retrieved material
\citep{ae1_lewis2020rag}, rather than a distinct kind of knowledge
or a synonym for skill. Retrieval can supply examples for structured
prediction \citep{ae1_rubin2022epr} or reusable reasoning templates
\citep{ae1_yang2024buffer}. In LEAN-LLM-OPT, reference retrieval
supports classification and demonstration construction, whereas
CSVQA retrieves and returns relevant information from the target
datasets. The latter supplies instance-specific information needed
to formulate the current model, rather than additional reference
examples.

Recent work describes agent skills as reusable procedural artifacts
organized around instructions, supporting resources, and conditions
for application \citep{ae1_zhou2026skillsurvey}. OpenAI's official
documentation similarly permits skills to contain reference material
and executable scripts, and to guide the use of search and fetch
tools \citep{ae1_openai2026skills,ae1_openai2026skilltools}.
Thus, a skill can incorporate knowledge and use retrieval during
execution. Reusability concerns the method, not an identical sequence
of actions or outputs across instances.

Viewed functionally, LEAN-LLM-OPT's type-tailored modeling procedures
are \emph{skill-like}: classification selects an applicable route;
reusable instructions and data-access routines organize how to
retrieve, interpret, and use information; and reference cases provide
supporting modeling knowledge. Here, problem types refer to
optimization structures, such as transportation and facility location.
The workflow generation agent integrates selected reference content
with procedural guidance to construct a demonstration. The model
generation agent then adapts this guidance to the target requirements,
uses the data tools to obtain the current instance's data, and builds
the formulation. For example, a transportation reference demonstrates
how supply, demand, and cost data enter a model, while the target
datasets provide the values for the new instance. For Others and
Mixture, the type-agnostic route instead guides construction through
an abstract model plan based on the current query and dataset schema.

We distinguish this procedural reuse from skill packaging and
learning. Our implementation expresses the procedures through
prompts, workflow-construction code, and data-access tools, rather
than standalone skill packages of the kind evaluated in SkillsBench
\citep{ae1_li2026skillsbench}. It also does not induce or update
skills from completed trajectories, as studied in Agent Workflow
Memory \citep{ae1_wang2025awm} and OptSkills
\citep{ae1_yang2026optskills}. The curated references and reusable
instructions remain fixed during evaluation, while demonstrations,
plans, and execution depend on the instance. The absence of skill
learning does not preclude skill-like procedural reuse; it delineates
the scope of the present framework.

Accordingly, LEAN-LLM-OPT remains an \emph{agentic workflow
construction framework} that combines reference modeling knowledge,
reusable modeling procedures, and access to external instance data.
The agents perform the reasoning and tool use through which these
resources are applied. This inference-time guidance can also
complement a fine-tuned model; our positioning does not imply that
it universally outperforms or replaces fine-tuning.

We have added \emph{Knowledge and Skills in LLM Agents} to
Section~1.2 (\emph{Related Works}) to explain the terminology and
relevant literature. We have also added \emph{Connecting Reference
Knowledge and Modeling Skills} to Section~3.3
(\emph{Novelties in Our Design and Key Messages}) to clarify the
roles of reference cases, procedural guidance, data-access tools,
and the two generation agents, together with the distinction between
procedural reuse and skill packaging or learning.
}
